% p1-03-symbolic-grammar

\section{P1 §3. Symbolic Grammar and Complexity Metric}

3. Symbolic Grammar and Complexity Metric

          Figure (fig-p1-ch3-grammar). Trinity grammar triptych --- basis abacus, open (p,m) lattice,
                                        five-spoke generator wheel.

  3.1 Algebraic Grammar Definition
  The symbolic grammar G$\varphi$ is defined over the alphabet

  A = {$\varphi$, pi, e, 3, Z},

  where $\varphi$ = (1+5)/2 is the golden ratio, pi approx 3.14159 is the circle constant, e approx
  2.71828 is Euler's number, and Z encodes integer-valued exponents and prefactors. The
  grammar generates expressions of the form

  F(theta) = n $\cdot$ 3k $\cdot$ $\varphi$p $\cdot$ pim $\cdot$ eq, 1

  where the parameter vector is

  theta := (n, k, p, m, q) $\in$ {1,…,9} $\times$ Z4.

  The admissible symbolic expression set at complexity level C $\in$ N is defined as

  S(C) := { F(theta) | theta $\in$ {1,…,9} $\times$ Z4, |theta|1 $\leq$ C }, 2

  where the ell1-norm complexity constraint is

  The prefactor n $\in$ {1,…,9} is treated as a single-digit integer whose absolute value is
  excluded from the complexity count; this reflects the convention that small integer
  prefactors represent negligible description cost relative to exponent choices. All
  F(theta) are positive real numbers by construction.

  3.2 Complexity Metric and Its Justification
  The ell1-norm complexity |theta|1 acts as a bounded description-length regularizer. Its
  adoption is motivated by three considerations:

  - Kolmogorov analogy. The number of bits required to specify an integer exponent tuple
  scales as log2 |theta|1 per entry, making the ell1 norm a conservative linear proxy for
  code length (Barron, Rissanen \& Yu, 1998; Chaitin, 1987).
  - Algebraic naturality. Expressions with small ell1 norm require the fewest
  multiplicative factors and are algebraically minimal in a precise sense.
  - Tractability. The induced hypothesis class HC has finite, analytically bounded
  cardinality (Section 3.4), enabling exhaustive enumeration for C $\leq$ 15.

  Sensitivity to this metric choice is discussed in Section 13.

  3.3 Logarithmic Embedding Structure
  Under the logarithmic map log : R^+ $\rightarrow$ R, expression (1) becomes

  log F(theta) = log n + k log 3 + p log $\varphi$ + m log pi + q log e. 3

  Defining the basis vector v := (log 3, log $\varphi$, log pi, 1) $\in$ R4, the image of S(C) in
  log-space is a bounded section of the rank-4 additive lattice

  log S(C) $\subset$ { log n + (k,p,m,q)$\cdot$v | (k,p,m,q)$\in$Z4, |(k,p,m,q)|1 $\leq$ C }. 4

  Algebraic closure of the Lucas ring Z[$\varphi$] (PhD: Ch.3 Trinity Identity) ensures stability
  under the arithmetic operations of the grammar. The non-degeneracy of the log-lattice
  is treated as a working independence assumption on {1, log 3, log $\varphi$, log pi} over Q, not
  as a theorem of the present paper. If a non-trivial rational relation among these
  generators were found, the hypothesis class would be quotient by that relation and all
  cardinalities and MDL penalties would need to be recomputed.

  3.4 Cardinality and Combinatorial Complexity

  Figure (fig-p1-ch3-search-space). Search-space triptych --- bounded exponent grid, |S(C)| growth,
                                          Bayesian prior.

  The cardinality of S(C) governs both computational cost and multiple-testing severity.
  Standard lattice-point counting in the ell1-ball in Z4 gives

  reflecting the volume of the 4-dimensional integer ell1-ball scaled by 9 prefactor
  choices. This quartic growth is significant: the hypothesis class grows rapidly but
  remains computationally enumerable for moderate complexity bounds. At C = 6, one
  finds |S(6)| = 11{,}601 expressions; at C = 10, |S(10)| = 75{,}249; at C = 15, |S(15)| =
  351{,}369. All of these are amenable to exhaustive enumeration on modern hardware.
  The effective number of independent tests in any significance assessment is Neff =
  |S(C)|; this sets the scale for all multiple-testing corrections in Section 6.

  The following table summarizes exact cardinalities at selected complexity levels:

  Note: exact counts include the 9 prefactor choices and the full signed ell1-ball volume
  formula.

  3.5 Representation Bias and the Induced Grammar Prior
  The grammar G$\varphi$ induces a non-uniform prior measure over R^+. The empirical
  counting measure

  PG$\varphi$(x) propto sumtheta $\in$ S(C) delta(x - F(theta)), 6

  smoothed with an exponential kernel of bandwidth sigma > 0, yields

  Psigma(x) = sumtheta $\in$ S(C) exp(-|x - F(theta)|/sigma). 7

  This distribution is concentrated near dense algebraic combinations of the basis
  elements and is generically non-uniform over R^+. The non-uniformity of Psigma must
  be explicitly accounted for in all null-model constructions (Section 5). Failure to
  account for this non-uniformity would systematically inflate apparent compressibility.

  3.6 Approximation Quality Functional
  Given a target constant X $\in$ R^+, the approximation quality at complexity C is
  measured by the minimum normalized absolute deviation:

  dC(X) := mintheta $\in$ S(C) (|F(theta) - X|)/(|X|). 8

  The function C ↦ dC(X) is monotonically non-increasing in C and converges to zero for
  sufficiently large C whenever X $\in$ R^+. The statistical question concerns the rate of
  this convergence relative to null expectations under properly specified baseline models.

\subsection{References}

- Vasilev, D. \textit{Flos Aureus Monograph: Trinity Constants}. DOI \href{https://zenodo.org/doi/10.5281/zenodo.19227877}{10.5281/zenodo.19227877}.
- Pellis, S. "The Fine-Structure Constant and the Golden Ratio." \href{https://vixra.org/pdf/2110.0117v4.pdf}{viXra 2110.0117 v4}.
- Sherbon, M. A. "Fundamental Nature of the Fine-Structure Constant." \href{https://hal.science/hal-02341850/document}{HAL hal-02341850}.
